\chapter{Capacity, Decay and Crowding}\label{ch:rs:capacity-decay-and-crowding}

From Monday 6 to Thursday 9 August 2007 the value factor of the US stock market lost 2.28 per cent and the momentum factor 3.43 per cent, ten and twelve times
their daily standard deviations, while the market itself rose 1.4 per cent. There was no news on the stocks. On Friday the value factor regained 1.42 per cent
and momentum 1.15 per cent. Khandani and Lo traced the episode to the unwinding of long--short equity books that many quantitative funds held in common, and to
market makers withdrawing their capital as it happened. A strategy's \omterm{def:rs:capacity-decay-and-crowding:capacity}{capacity} is not a property of its signal alone: it depends on its own costs, which grow
with its size, and on everyone else trading the same names. This chapter draws \omterm{def:rs:capacity-decay-and-crowding:capacity}{capacity curves} with \texttt{firm.capacity}, looks at alpha decay in production,
and simulates an \omterm{def:rs:capacity-decay-and-crowding:crowding}{unwind}.

\section{Capacity curves}

\begin{definition}[Strategy capacity, capacity curve, profit-maximising size]\label{def:rs:capacity-decay-and-crowding:capacity}
A \emph{capacity curve}\index{capacity curve} gives a strategy's net return, net Sharpe ratio and dollar profit against the size of the capital it runs, each
size traded with its costs. The \emph{profit-maximising size}\index{profit-maximising size} is the size at which dollar profit peaks. A strategy's
\emph{capacity}\index{strategy capacity} is the size beyond which a stated criterion fails: here, the size at which its net Sharpe ratio falls to half its best.
\end{definition}

\begin{proposition}[Profit and size under concave costs]\label{prop:rs:capacity-decay-and-crowding:size}
If a strategy's gross return $g$ does not depend on its size $A$ and its costs, as a fraction of capital, grow as $k\sqrt A$ (trades proportional to capital,
square-root impact), its dollar profit $A(g - k\sqrt A)$ is maximised at $A^* = (2g/3k)^2$, where its net return is $g/3$.
\end{proposition}

\begin{proof}
The derivative $g - \frac32 k\sqrt A$ vanishes at $\sqrt{A^*} = 2g/3k$, and the net return there is $g - k\sqrt{A^*} = g/3$.
\end{proof}

The \omterm{def:rs:capacity-decay-and-crowding:capacity}{profit-maximising size} is therefore larger than the size that maximises the return on capital, which is the smallest; Berk and Green built a theory of
active management on this decreasing return to scale. \Cref{fig:rs:capacity-decay-and-crowding:capacity} draws chapter~27's fifty-name book at ten sizes from
\$10 million to \$30 billion, in two versions. The fixed rule (costs ignored, the forecast smoothed over fifty days) trades the same way at every size: its net
Sharpe ratio falls from 1.79 at \$10 million to 0.34 at \$5 billion and $-0.28$ at \$10 billion, halving at \$2.02 billion, and its dollar profit peaks at \$3
billion, \$146 million a year. The book re-optimised at each size, with the costs inside and a one-day smoothing, peaks at a Sharpe ratio of 2.43 at \$1 billion
and is still at 1.36 at \$30 billion, where its profit, \$806 million a year, is still growing: re-optimising at size is the largest single extension of
\omterm{def:rs:capacity-decay-and-crowding:capacity}{capacity}, because it changes how much a book trades as it grows.

\begin{omfigure}
\begin{tikzpicture}
\begin{semilogxaxis}[name=a,width=6.8cm,height=5.4cm,xlabel={\small fund size (\$)},ylabel={\small net Sharpe ratio},tick label style={font=\footnotesize},
  xmin=7e6,xmax=4e10,ymin=-2.1,ymax=2.8,grid=major,grid style={gray!25},legend pos=south west,legend style={font=\scriptsize},table/col sep=comma]
\addplot[black!60,thick,mark=o] table[x=size,y=sr_fixed]{figdata/research/28-capacity-decay-and-crowding/capacity.csv};
\addplot[omDef,very thick,mark=*] table[x=size,y=sr_opt]{figdata/research/28-capacity-decay-and-crowding/capacity.csv};
\legend{fixed rule,re-optimised}
\end{semilogxaxis}
\begin{semilogxaxis}[at={(a.east)},anchor=west,xshift=1.3cm,width=6.8cm,height=5.4cm,xlabel={\small fund size (\$)},ylabel={\small profit (\$ million a year)},
  tick label style={font=\footnotesize},xmin=7e6,xmax=4e10,ymin=-320,ymax=850,grid=major,grid style={gray!25},table/col sep=comma]
\addplot[black!60,thick,mark=o] table[x=size,y=profit_fixed]{figdata/research/28-capacity-decay-and-crowding/capacity.csv};
\addplot[omDef,very thick,mark=*] table[x=size,y=profit_opt]{figdata/research/28-capacity-decay-and-crowding/capacity.csv};
\end{semilogxaxis}
\end{tikzpicture}

\omcaption{\omterm{def:rs:capacity-decay-and-crowding:capacity}{Capacity curves} of chapter~27's fifty-name book: net Sharpe ratio (left) and dollar profit (right) against fund size, for the fixed rule and the book
re-optimised at each size (profit below $-\$300$ million is drawn at $-300$). Data: \texttt{rs\_capacity.curve}.}\label{fig:rs:capacity-decay-and-crowding:capacity}
\end{omfigure}

\section{Alpha decay in production}

A strategy's \omterm{def:rs:capacity-decay-and-crowding:capacity}{capacity} changes after it goes live, for reasons chapter~13 measured: the signal decays as others find it, and the post-publication decline of
published factors is the public record of that. \omterm{def:rs:capacity-decay-and-crowding:capacity}{Capacity} decays twice over, since the gross return falls and the costs, now shared with the competitors trading
the same names, rise. The \omterm{def:rs:capacity-decay-and-crowding:capacity}{capacity curve} should therefore be redrawn with the live \omterm{def:rs:anatomy-of-a-predictor:ic}{information coefficient} and with the costs measured on the firm's own fills
(chapters~23 and~27), and a strategy should be sized below its estimated \omterm{def:rs:capacity-decay-and-crowding:capacity}{profit-maximising size}, where the curve is flat: the profit lost by being smaller is
small, the protection against a curve that has moved is large.

\section{Detecting crowding}

\begin{definition}[Crowding, crowded trade, comomentum, unwind]\label{def:rs:capacity-decay-and-crowding:crowding}
\emph{Crowding}\index{crowding} is the presence of many investors in the same positions; a \emph{crowded trade}\index{crowded trade} is a position held by enough
of them that their joint exit would move prices. \emph{Comomentum}\index{comomentum} (Lou and Polk) is the average correlation of the abnormal returns of the
stocks a momentum strategy holds, a measure of how many traders act on them. An \emph{unwind}\index{unwind} is the rapid reduction of a crowded position by some
of its holders, which moves the prices at which the others mark theirs.
\end{definition}

Stein named the problem: for a strategy whose value has no anchor in fundamentals, an arbitrageur cannot know how many of its peers are entering the same trade,
and the leverage each chooses in private creates a fire-sale risk for all. Pedersen analysed why traders crowd in and why they run. Three monitors make \omterm{def:rs:capacity-decay-and-crowding:crowding}{crowding}
visible (\cref{lst:rs:capacity-decay-and-crowding:curve}): the correlation among the returns of a strategy's own names, net of the risk model's factors (\omterm{def:rs:capacity-decay-and-crowding:crowding}{comomentum}
for a momentum book: high values predicted reversals); the overlap of the firm's book with estimates of others' books (from regulatory holdings filings, or from
the returns of competitors' products); and short interest and days to cover on the short side (Book~1, chapter~16). None is conclusive, and each is noisy; they
belong on the risk dashboard next to the book's liquidity.

\section{August 2007}

The chapter's \omterm{def:rs:capacity-decay-and-crowding:crowding}{unwind} (\cref{lst:rs:capacity-decay-and-crowding:unwind}) puts five funds of \$2 billion each in long--short books of 500 stocks, four times their
capital in gross exposure, a hundred names a side; each stock trades \$200 million a day with a daily volatility of 2\%. Each fund ranks the stocks on a mixture of
a common score and its own, with a weight that sets how much their books overlap. One fund sells a fraction of its book over three days. Each day its trades
move prices by a square-root impact, 30\% of which stays while the rest decays with a half-life of a day, and the other four funds mark their books to the moved
prices.

\begin{center}
\small
\begin{tabular}{@{}lrrrr@{}}
\toprule
 & \multicolumn{4}{c}{overlap of the books (cosine)}\\
\cmidrule(l){2-5}
fraction of the book sold & 0.04 & 0.18 & 0.43 & 1.00\\
\midrule
10\% & $-0.03\%$ & $-0.18\%$ & $-0.42\%$ & $-0.97\%$\\
25\% & $-0.05\%$ & $-0.28\%$ & $-0.66\%$ & $-1.54\%$\\
50\% & $-0.08\%$ & $-0.40\%$ & $-0.93\%$ & $-2.17\%$\\
100\% & $-0.11\%$ & $-0.56\%$ & $-1.32\%$ & $-3.07\%$\\
\bottomrule
\end{tabular}
\end{center}

The table gives the other funds' peak loss, as a fraction of their capital; it comes on the third day, the last of the selling, in every case. The loss grows
in proportion to the overlap and with the square root of the fraction sold, since impact is a square root: selling a quarter of the book costs the others half as
much as selling all of it. With identical books and the whole book sold, the others lose 3.07\% of their capital without trading, the seller 2.00\% (it also pays
for its trades). Then the \omterm{def:rs:transaction-costs-in-research:impact}{temporary impact} decays: by day fifteen the others' loss is 1.30\%, the permanent share (\cref{fig:rs:capacity-decay-and-crowding:unwind}).
The recovery is the Friday of August 2007, and the ones who lost most were those who sold at the bottom, locking in the temporary part. A fund that knows the
loss is temporary and has the capital to hold on loses only the permanent part; a fund with a stop-loss becomes the next seller, which is how an \omterm{def:rs:capacity-decay-and-crowding:crowding}{unwind} spreads.

\begin{omfigure}
\begin{tikzpicture}
\begin{axis}[width=11cm,height=5.8cm,xlabel={days from the start of the selling},ylabel={loss of the other funds (\% of capital)},tick label style={font=\small},
  label style={font=\small},xmin=0,xmax=15,ymin=-3.3,ymax=0.2,grid=major,grid style={gray!25},legend columns=4,
  legend style={font=\footnotesize,at={(0.5,-0.25)},anchor=north,draw=none,fill=none,/tikz/every even column/.append style={column sep=8pt}},
  table/col sep=comma]
\addplot[omDef!35,thick] table[x=day,y=f10]{figdata/research/28-capacity-decay-and-crowding/unwind_paths.csv};
\addplot[omDef!60,thick] table[x=day,y=f25]{figdata/research/28-capacity-decay-and-crowding/unwind_paths.csv};
\addplot[omDef!85,thick] table[x=day,y=f50]{figdata/research/28-capacity-decay-and-crowding/unwind_paths.csv};
\addplot[omThm,very thick] table[x=day,y=f100]{figdata/research/28-capacity-decay-and-crowding/unwind_paths.csv};
\draw[black!50,densely dotted] (axis cs:3,-3.3) -- (axis cs:3,0.2);
\legend{10\% sold,25\% sold,50\% sold,all sold}
\end{axis}
\end{tikzpicture}

\omcaption{The average loss of the four funds that did not sell, with identical books, when one fund sells part of its book over three days (the dotted line
marks the last day of selling). The permanent share of the impact is what remains. Data: \texttt{rs\_capacity.unwind}.}\label{fig:rs:capacity-decay-and-crowding:unwind}
\end{omfigure}

\section{Tutorial: the week of 6 August 2007}\label{tut:rs:capacity-decay-and-crowding:tut}

\begin{tutorial}
\textbf{Goal.} Draw the \omterm{def:rs:capacity-decay-and-crowding:capacity}{capacity curves} of chapter~27's book, find the \omterm{def:rs:capacity-decay-and-crowding:capacity}{profit-maximising size} and the size at which the net Sharpe ratio halves, and simulate
the \omterm{def:rs:capacity-decay-and-crowding:crowding}{unwind} over a grid of overlaps and fractions sold. \textbf{End state:} the table, \cref{fig:rs:capacity-decay-and-crowding:capacity,fig:rs:capacity-decay-and-crowding:unwind},
and the \omterm{def:rs:risk-models:families}{factor returns} of August 2007 from \texttt{data/research/ff\_aug2007.csv}.
\begin{tutsteps}
\item \textbf{\omterm{def:rs:capacity-decay-and-crowding:capacity}{Capacity} helpers}: the curve, its peak and the size at a fraction of the best.
\omcode{code/firm/capacity/firm_capacity.py}{29}{46}{Capacity curve, profit-maximising size and halving size.}\label{lst:rs:capacity-decay-and-crowding:curve}
\item \textbf{The \omterm{def:rs:capacity-decay-and-crowding:crowding}{unwind}}: trades at the open, square-root impact, permanent and decaying parts, marking.
\omcode{code/firm/capacity/firm_capacity.py}{64}{89}{Funds marking overlapping books while one sells.}\label{lst:rs:capacity-decay-and-crowding:unwind}
\item \textbf{Run} \texttt{rs\_capacity.curve}, \texttt{summary}, \texttt{unwind(theta, fraction)}, \texttt{fig\_capacity.py}, and \texttt{rs\_fetch\_aug2007.py} once.
\end{tutsteps}
\textbf{What to change next.} Give the non-selling funds a stop-loss at 1\% and let them sell when it is hit: the cascade; make the permanent share depend on
whether market makers stay (Khandani and Lo's withdrawal of capital).
\end{tutorial}

\section{Build: capacity and crowding}\label{bld:rs:capacity-decay-and-crowding:capacity}

\begin{build}
\bfield{Purpose} Every strategy is sized from its \omterm{def:rs:capacity-decay-and-crowding:capacity}{capacity curve} and monitored for \omterm{def:rs:capacity-decay-and-crowding:crowding}{crowding}; the risk of an \omterm{def:rs:capacity-decay-and-crowding:crowding}{unwind} by others is a scenario in the firm's stress
tests.
\bfield{Interface} \texttt{capacity\_curve(sizes, net\_returns, vols)}, \texttt{profit\_maximising(curve)}, \texttt{size\_at\_fraction(curve, frac, key)},
\texttt{avg\_pairwise\_corr(R)}, \texttt{comomentum(R\_long, R\_short)}, \texttt{overlap(w1, w2)}, \texttt{unwind(books, capitals, seller, fraction, days, adv,
sigma, eta, permanent, half\_life, horizon)}.
\bfield{Rules} \omterm{def:rs:capacity-decay-and-crowding:capacity}{Capacity curves} are drawn from \omterm{def:rs:vectorised-backtests:backtest}{backtests} re-run at each size with the firm's fitted costs, and redrawn with live \omterm{def:rs:anatomy-of-a-predictor:ic}{information coefficients}; a
strategy runs below its \omterm{def:rs:capacity-decay-and-crowding:capacity}{profit-maximising size}; \omterm{def:rs:capacity-decay-and-crowding:crowding}{crowding} monitors are reviewed with the book's liquidity.
\bfield{Acceptance tests} \texttt{code/firm/capacity/tests/}: the profit-maximising and halving sizes of a strategy with square-root costs against their closed
forms; pairwise correlation and \omterm{def:rs:capacity-decay-and-crowding:crowding}{comomentum} of a planted common factor; overlap by hand; an \omterm{def:rs:capacity-decay-and-crowding:crowding}{unwind} in which the same book loses and the opposite book gains the
same, \omterm{def:rs:transaction-costs-in-research:impact}{temporary impact} that comes back, and \omterm{def:rs:transaction-costs-in-research:impact}{permanent impact} that does not.
\bfield{Stretch} Cascades with stop-losses; \omterm{def:rs:capacity-decay-and-crowding:crowding}{crowding} estimated from regulatory holdings filings; \omterm{def:rs:capacity-decay-and-crowding:capacity}{capacity} shared among the firm's strategies.
\end{build}

\begin{omsources}
\item A.~E.~Khandani and A.~W.~Lo, ``What happened to the quants in August 2007? Evidence from factors and transactions data'', \emph{Journal of Financial Markets}
14(1), 2011.
\item D.~Lou and C.~Polk, ``Comomentum: inferring arbitrage activity from return correlations'', \emph{Review of Financial Studies} 35(7), 2022.
\item L.~H.~Pedersen, ``When everyone runs for the exit'', NBER Working Paper 15297, 2009.
\item J.~C.~Stein, ``Presidential address: sophisticated investors and market efficiency'', \emph{Journal of Finance} 64(4), 2009.
\item J.~B.~Berk and R.~C.~Green, ``Mutual fund flows and performance in rational markets'', \emph{Journal of Political Economy} 112(6), 2004.
\item K.~R.~French, Data Library, Fama/French 3 Factors and Momentum Factor (daily).
\end{omsources}

\section{Exercises}

\begin{exercise}[$\star$]\label{exo:rs:capacity-decay-and-crowding:1}
A strategy earns 10\% a year gross and its costs reach 10\% at \$10 billion, growing as the square root of size. Find its \omterm{def:rs:capacity-decay-and-crowding:capacity}{profit-maximising size} and its net return
there.
\end{exercise}

\begin{exercise}[$\star$]\label{exo:rs:capacity-decay-and-crowding:2}
Why is the \omterm{def:rs:capacity-decay-and-crowding:capacity}{profit-maximising size} larger than the size that maximises the return on capital?
\end{exercise}

\begin{exercise}[$\star$]\label{exo:rs:capacity-decay-and-crowding:3}
In the chapter's \omterm{def:rs:capacity-decay-and-crowding:crowding}{unwind} with identical books, how large is one day's price push on a long name when the whole book is sold over three days?
\end{exercise}

\begin{exercise}[$\star\star$]\label{exo:rs:capacity-decay-and-crowding:4}
Why does the other funds' loss grow with the square root of the fraction sold, and in proportion to the overlap?
\end{exercise}

\begin{exercise}[$\star\star$]\label{exo:rs:capacity-decay-and-crowding:5}
In the week of 6 August 2007 the market rose while value and momentum fell by ten standard deviations. What does that say about who was trading?
\end{exercise}

\begin{exercise}[$\star\star$]\label{exo:rs:capacity-decay-and-crowding:6}
Why does re-optimising at each size extend \omterm{def:rs:capacity-decay-and-crowding:capacity}{capacity} so much more than any change to the signal?
\end{exercise}

\begin{exercise}[$\star\star\star$]\label{exo:rs:capacity-decay-and-crowding:7}
\emph{Coding.} With \texttt{rs\_capacity.unwind}, compare the other funds' loss at day fifteen with their peak loss, and relate the ratio to the permanent share.
\end{exercise}

\begin{exercise}[$\star\star\star$]\label{exo:rs:capacity-decay-and-crowding:8}
\emph{Find the flaw.} ``Our strategy's \omterm{def:rs:capacity-decay-and-crowding:capacity}{capacity} is \$5 billion: at that size its \omterm{def:rs:vectorised-backtests:backtest}{backtest} still has a Sharpe ratio of 0.34 after costs.''
\end{exercise}

\section{Problem: The Week of 6 August 2007}

\begin{problem}[{Weekend problem --- how big, and with whom}]\label{pb:rs:capacity-decay-and-crowding:1}
Chapter~27's book at many sizes, and the chapter's \omterm{def:rs:capacity-decay-and-crowding:crowding}{unwind}.

\medskip\noindent\textbf{Part I --- \omterm{def:rs:capacity-decay-and-crowding:capacity}{Capacity}.}
\begin{enumerate}
\item State and prove the \omterm{def:rs:capacity-decay-and-crowding:capacity}{profit-maximising size} under square-root costs.
\item Give the fixed rule's net Sharpe ratio from \$10 million to \$30 billion.
\item Where does its profit peak, and where does its Sharpe ratio halve?
\item What does re-optimising at each size change?
\item How would you size the strategy, and why below its \omterm{def:rs:capacity-decay-and-crowding:capacity}{profit-maximising size}?
\end{enumerate}

\medskip\noindent\textbf{Part II --- Decay and \omterm{def:rs:capacity-decay-and-crowding:crowding}{crowding}.}
\begin{enumerate}[resume]
\item Why does \omterm{def:rs:capacity-decay-and-crowding:capacity}{capacity} decay after a strategy goes live?
\item What are Stein's two complications?
\item What is \omterm{def:rs:capacity-decay-and-crowding:crowding}{comomentum}, and what did high values predict?
\item Which monitors would you put on the dashboard?
\end{enumerate}

\medskip\noindent\textbf{Part III --- The \omterm{def:rs:capacity-decay-and-crowding:crowding}{unwind}.}
\begin{enumerate}[resume]
\item Describe the simulated funds and the impact model.
\item Give the other funds' peak losses over the grid.
\item When does the peak come, and what remains at day fifteen?
\item What did the value and momentum factors do from 6 to 10 August 2007?
\item What did Khandani and Lo identify?
\end{enumerate}

\medskip\noindent\textbf{Part IV --- The verdict.}
\begin{enumerate}[resume]
\item State the \emph{named result}: the peak loss of the non-selling funds against the fraction of the overlapping book sold, and the size at which the strategy's
net Sharpe ratio halves.
\item What should a fund do on the third day of an \omterm{def:rs:capacity-decay-and-crowding:crowding}{unwind} it did not start?
\item How would stop-losses change the table?
\item How does the firm's own set of strategies crowd itself (chapter~27's netting)?
\item What would you ask a new strategy's author about \omterm{def:rs:capacity-decay-and-crowding:capacity}{capacity}?
\item In one sentence: what determines a strategy's \omterm{def:rs:capacity-decay-and-crowding:capacity}{capacity}?
\end{enumerate}
\end{problem}

\section{Interview questions}

\begin{interviewq}[$\star$ \iqroles{researcher, risk}]\label{iq:rs:capacity-decay-and-crowding:1}
What happened to quantitative equity funds in August 2007?
\end{interviewq}

\begin{interviewq}[$\star\star$ \iqroles{researcher}]\label{iq:rs:capacity-decay-and-crowding:2}
How do you estimate a strategy's \omterm{def:rs:capacity-decay-and-crowding:capacity}{capacity}?
\end{interviewq}

\begin{interviewq}[$\star\star$ \iqroles{researcher, risk}]\label{iq:rs:capacity-decay-and-crowding:3}
How would you tell whether your strategy is crowded?
\end{interviewq}

\begin{interviewq}[$\star\star$ \iqroles{risk}]\label{iq:rs:capacity-decay-and-crowding:4}
Another fund with a book like yours is liquidating. What do you do?
\end{interviewq}

\begin{interviewq}[$\star\star$ \iqroles{researcher, trader}]\label{iq:rs:capacity-decay-and-crowding:5}
Why can a strategy's dollar profit keep rising after its Sharpe ratio has started to fall?
\end{interviewq}

\begin{interviewq}[$\star\star\star$ \iqroles{researcher}]\label{iq:rs:capacity-decay-and-crowding:6}
Model the loss of a fund that holds a book overlapping a seller's, with square-root impact and a permanent share, and derive how it scales with the fraction sold.
\end{interviewq}
